\begin{proof}
\emph{Setup.} The cell structure of $S^{5}$, its boundary formulas $\partial e_{2r+1} = (g - 1)e_{2r}$, $\partial e_{2r} = Ne_{2r-1}$ and the diagonal $\Delta_{0}$ are those of clauses (2) and (3) of Theorem \claimref{synthesis.cohomology-of-the-quarter-turn-lens-spaces-from-the-equivariant-cell-structure}; the same construction for the half turn $T = -1$ on the two planes of $\HH$ gives the complex $D$ of $S^{3}$ (the formulas of that theorem for the group of order two, which its check covers). In both, the closed cells are balls: the closure of an even cell is a closed disc, and the closure of an odd cell is the join of a sphere with a closed arc; the closure of $g^{a}e_{2r+1}$ consists of that cell, its two ends $g^{a}e_{2r}$, $g^{a+1}e_{2r}$ and the sphere $S^{2r-1}$, and the closure of $g^{a}e_{2r}$ of that cell and $S^{2r-1}$. The quarter turn is the rotor action $J(v) = \tilde U v U/8$ with $U = (1 + e_{1}e_{2})(1 + e_{3}e_{4})(1 + e_{5}e_{6})$, and the even subalgebra of $\Cl(3,0)$ is spanned by $1$ and three anticommuting bivectors of square $-1$, with $q\tilde q$ a multiplicative positive definite scalar form (check). The group acts freely on $S^{5} \times S^{3}$ because it does on the first factor, and it permutes the product cells $g^{a}e_{k} \times T^{b}d_{q}$ freely; there are $6 \cdot 4 \cdot 2 = 48$ orbits, with representatives $e_{k} \otimes T^{b}d_{q}$, and an invariant cochain is a function on them.

\emph{(1).} Given $F$, the point $[v, F(v)/|F(v)|]$ of $M$ depends only on the class of $v$, because $(J(v), F(J(v))) = (J(v), -F(v)) = g \cdot (v, F(v))$; so $s[v] = [v, F(v)/|F(v)|]$ is a continuous section. Conversely, for a section $s$ and $v \in S^{5}$ the orbit $s[v]$ contains exactly one point with first coordinate $v$ (the action is free on $S^{5}$); calling it $(v, F(v))$ defines $F$ with $F(J(v)) = -F(v)$, continuous because $S^{5} \times S^{3} \to M$ is a covering. Only the first direction is used below.

\emph{(2).} On each factor the identity $\partial\Delta_{1}(e_{k}) + \Delta_{1}(\partial e_{k}) = \Delta_{0}(e_{k}) + t\Delta_{0}(e_{k})$ modulo $2$ is a finite statement about the listed terms, with $\Delta_{1}$ extended to $g^{a}e_{k}$ by the diagonal action; the check verifies it for $k \leq 5$ on $C$ and $k \leq 3$ on $D$, together with the carrier property (every cell occurring in $\Delta_{0}(e_{k})$ or $\Delta_{1}(e_{k})$ lies in the closure of $e_{k}$ described above). For example $\Delta_{0}(e_{1}) + t\Delta_{0}(e_{1}) = e_{0} \otimes e_{1} + e_{1} \otimes ge_{0} + e_{1} \otimes e_{0} + ge_{0} \otimes e_{1}$, which is $\partial(e_{1} \otimes e_{1})$ modulo $2$. On the product, modulo $2$ and with the middle factors interchanged (which is an isomorphism of complexes modulo $2$ under which the exchange $t$ of $P \otimes P$ becomes $t \otimes t$), the commutator with the boundary of $\Delta_{1}^{C} \otimes t\Delta_{0}^{D} + \Delta_{0}^{C} \otimes \Delta_{1}^{D}$ is
\[
(\Delta_{0}^{C} + t\Delta_{0}^{C}) \otimes t\Delta_{0}^{D} + \Delta_{0}^{C} \otimes (\Delta_{0}^{D} + t\Delta_{0}^{D}) = t\Delta_{0}^{C} \otimes t\Delta_{0}^{D} + \Delta_{0}^{C} \otimes \Delta_{0}^{D},
\]
since $\Delta_{0}^{C}$, $\Delta_{0}^{D}$ and their transposes are chain maps. The factors on $C$ are equivariant for $g$ and those on $D$ for $T$, so the product maps are equivariant for $g \otimes T$; the closure of a product cell is the product of the closures, so the product maps are carried by closed cells. The check verifies the identity and the carrier property directly on all $48$ orbit representatives, for this formula and for $\Delta_{1}^{C} \otimes \Delta_{0}^{D} + t\Delta_{0}^{C} \otimes \Delta_{1}^{D}$, whose commutator with the boundary is the same by the same computation.

\emph{(3).} Write a cochain on the representatives as $\varphi_{k,q} = (\varphi(e_{k} \otimes d_{q}), \varphi(e_{k} \otimes Td_{q}))$; invariance gives $\varphi(g^{a}e_{k} \otimes T^{b}d_{q}) = \varphi(e_{k} \otimes T^{b-a}d_{q})$. The cocycle condition in degree $3$ is the vanishing on the boundaries of the cells of dimension $4$:
\[
e_{1} \otimes T^{b}d_{3}: \ \varphi_{0,3}(b{+}1) - \varphi_{0,3}(b) = \varphi_{1,2}(b{+}1) - \varphi_{1,2}(b); \qquad
e_{2} \otimes T^{b}d_{2}: \ 2\,(\varphi_{1,2}(0) + \varphi_{1,2}(1)) + \varphi_{2,1}(0) + \varphi_{2,1}(1) = 0;
\]
\[
e_{3} \otimes T^{b}d_{1}: \ \varphi_{2,1}(b{+}1) - \varphi_{2,1}(b) = \varphi_{3,0}(b{+}1) - \varphi_{3,0}(b); \qquad
e_{4} \otimes T^{b}d_{0}: \ 2\,(\varphi_{3,0}(0) + \varphi_{3,0}(1)) = 0,
\]
where $2 = 4/2$ counts the translates $g^{a}e_{k}$ with $a$ of each parity in $Ne_{k}$. The cochain $\theta$, $\theta_{0,3} = \theta_{1,2} = (1, 0)$, $\theta_{2,1} = (-1, -1)$, $\theta_{3,0} = (0, 0)$, satisfies them; its value on the cycle $(1 + T)(e_{0} \otimes d_{3})$, the fibre sphere over the base point, is $1$. Its reduction $a$ is $1$ on the cells $g^{a}e_{0} \otimes T^{a}d_{3}$, $g^{a}e_{1} \otimes T^{a}d_{2}$ and on all cells $g^{a}e_{2} \otimes T^{b}d_{1}$. The cochain $a \cup_{1} a$ is evaluated on the eight representatives of dimension $5$ in the check: it is $1$ on $e_{4} \otimes d_{1}$ and $e_{4} \otimes Td_{1}$ and $0$ on $e_{5} \otimes T^{b}d_{0}$, $e_{3} \otimes T^{b}d_{2}$, $e_{2} \otimes T^{b}d_{3}$. The surviving term is transparent: on $e_{4} \otimes d_{1}$ the part $\Delta_{1}^{C}(e_{4}) \otimes t\Delta_{0}^{D}(d_{1})$ pairs a cell of base dimension $3$ with one of base dimension $2$ and total fibre dimension $1$, on which $a \otimes a$ vanishes, and in $\Delta_{0}^{C}(e_{4}) \otimes \Delta_{1}^{D}(d_{1})$, with $\Delta_{1}^{D}(d_{1}) = d_{1} \otimes d_{1}$, only the term $e_{2} \otimes e_{2}$ of $\Delta_{0}^{C}(e_{4})$ gives two cells in the support of $a$, namely $(e_{2} \otimes d_{1}) \otimes (e_{2} \otimes d_{1})$, with value $\theta_{2,1}(0)^{2} = 1$. On the other representatives the terms cancel in pairs. The cochain $\pi^{*}\varepsilon_{5}$ is $1$ on $e_{5} \otimes T^{b}d_{0}$ and $0$ elsewhere. For $\beta$: $\delta\beta(e_{5} \otimes T^{b}d_{0}) = \beta((g - 1)e_{4} \otimes T^{b}d_{0}) = \beta(e_{4} \otimes T^{b+1}d_{0}) - \beta(e_{4} \otimes T^{b}d_{0}) = \pm 1$ and $\delta\beta(e_{4} \otimes T^{b}d_{1}) = \beta(e_{4} \otimes (T - 1)T^{b}d_{0}) = \pm 1$, and $\delta\beta$ vanishes on the other cells of dimension $5$, none of whose boundaries contains a cell $e_{4} \otimes T^{b}d_{0}$ with odd multiplicity. So $a \cup_{1} a + \pi^{*}\varepsilon_{5} = \delta\beta$ modulo $2$. That neither side is a coboundary, and that the other three choices of higher diagonal give the same cochain, is verified in the check by elimination over the field of two elements on the eight coboundaries of the basis cochains of degree $4$.

\emph{(4).} Suppose $(\sigma, h_{0}, h_{1})$ exist. The cochain $\sigma^{*}\theta$ is an invariant integral cocycle of degree $3$ of $C$, so $\sigma^{*}\theta = m\,\varepsilon_{3}$ with $0 = \delta(m\varepsilon_{3}) = 4m\,\varepsilon_{4}$ (cited theorem, clause (3)): $\sigma^{*}\theta = 0$, and $\sigma^{*}a = 0$. Apply $a \otimes a$ to the identity for $h_{1}$, evaluated on a chain $c$ of degree $5$, modulo $2$. Since $a \otimes a$ is a cocycle of $P \otimes P$, $(a \otimes a)(\partial h_{1}c) = 0$; $(a \otimes a)(h_{1}\partial c) = \delta\gamma(c)$ with $\gamma := (a \otimes a)h_{1}$, an invariant cochain of $C$ of degree $4$; $(a \otimes a)\Delta_{1}^{P}\sigma = \sigma^{*}(a \cup_{1} a)$; $(a \otimes a)(\sigma \otimes \sigma)\Delta_{1}^{C} = (\sigma^{*}a \otimes \sigma^{*}a)\Delta_{1}^{C} = 0$; and $(a \otimes a)(h_{0} + th_{0}) = 2(a \otimes a)h_{0} = 0$ because $a \otimes a$ is symmetric. Hence $\sigma^{*}(a \cup_{1} a) = \delta\gamma$ modulo $2$. By (3) and $\pi\sigma = 1$, $\sigma^{*}(a \cup_{1} a) = \varepsilon_{5} + \delta\sigma^{*}\beta$. So $\varepsilon_{5}$ is the coboundary modulo $2$ of an invariant cochain of $C$ of degree $4$, a multiple of $\varepsilon_{4}$; but $\delta\varepsilon_{4} = 0$. This contradiction proves the first statement. For the second: the equivariant map $\phi : C \to D$ with $\phi(e_{0}) = d_{0}$, $\phi(e_{1}) = d_{1}$, $\phi(e_{2}) = 2d_{2}$, $\phi(e_{3}) = (1 - T)d_{3}$, $\phi(e_{4}) = \phi(e_{5}) = 0$ and $\phi(gc) = T\phi(c)$ is an integral chain map (check), and $\sigma = (1 \otimes \phi)\Delta_{0}^{C}$ is then an equivariant chain map $C \to P$ whose projection keeps the terms $e_{k} \otimes (\text{vertex})$ of $\Delta_{0}^{C}(e_{k})$: $\pi\sigma = 1$.

\emph{(5).} $Y = S^{5} \times S^{3}$ with the product cells is a finite regular cell complex (products of closed balls), with cellular chains $P$ by (P), and the action is free and cellular; by (G) the cochains of $M$ are the invariant cochains of $P$. By (2) the maps $\Delta_{0}^{P}$, $\Delta_{1}^{P}$ meet the hypotheses of (S), which is clause (iii) of Fact \claimref{synthesis.imported-steenrod-cup-one-square-is-a-natural-cohomology-operation} ($\Delta_{0}^{P}$ is a chain map preserving the augmentation because its factors are). So, by (S) and (3), $Sq^{2}(\rho[\theta]) = [a \cup_{1} a] = [\pi^{*}\varepsilon_{5}] = p^{*}[\varepsilon_{5}] = p^{*}(xu^{2})$, the last two equalities by (P) and by clause (3) of the cited theorem. Let $s$ be a continuous section. Then $s^{*}[\theta] \in H^{3}(\Sig_{2}; \Z) = 0$ (cited theorem, clause (3)), and by naturality of the reduction modulo $2$ and of $Sq^{2}$ and by $p \circ s = 1$,
\[
xu^{2} = s^{*}p^{*}(xu^{2}) = s^{*}Sq^{2}(\rho[\theta]) = Sq^{2}(\rho\,s^{*}[\theta]) = Sq^{2}(0) = 0,
\]
where $Sq^{2}(0) = 0$ by (S) on $\Sig_{2}$ with the zero cocycle. But $xu^{2} \neq 0$ (cited theorem, clause (3)). So $p$ has no section and, by (1), no $F$ exists. A map $F$ on $S^{2n+1}$, $n \geq 2$, odd under the quarter turn, restricts to the unit sphere of the first three planes, which the quarter turn preserves.

\emph{(6).} For the group of order $d$, $4 \mid d$, the cocycle equations of (3) hold with $d/2$ in place of $2$, and $\theta_{d}$ with $\theta_{2,1} = (-d/4, -d/4)$ solves them; the surviving term of the square has value $(d/4)^{2}$, and the computation of $\delta\beta$ is unchanged. The check runs $d = 4, 8, 12, 16$ with the factor identity of (2) verified for each order. For $d \equiv 2$ modulo $4$ the second equation gives $\varphi_{2,1}(0) + \varphi_{2,1}(1) = -(d/2)(\varphi_{1,2}(0) + \varphi_{1,2}(1))$ and the third and fourth give $\varphi_{2,1}(1) - \varphi_{2,1}(0) = -2\varphi_{3,0}(0)$, so $(d/2)(\varphi_{1,2}(0) + \varphi_{1,2}(1))$ is even and $\varphi_{1,2}(0) + \varphi_{1,2}(1)$ is even; by the first equation $\varphi_{0,3}(0) + \varphi_{0,3}(1) \equiv \varphi_{1,2}(0) + \varphi_{1,2}(1)$ modulo $2$, so the value on the fibre sphere is even. The remaining controls are the check's.
\end{proof}
